\cvsection{Expériences professionnelles}

\begin{cventries}
  \cventry
    {Stagiaire ingénieur de recherche en IA}
    {INRIA}
    {Bordeaux, France}
    {Mars 2026 -- Septembre 2026}
    {
      \begin{cvitems}
        \item {Développement en autonomie d’une application full-stack intégrant un agent IA pour aider les chercheurs à retrouver des articles scientifiques.}
        \item {Orchestration de l’agent multi-étapes sans framework dédié, avec l’API OpenAI, des prompts adaptés au cas d’usage et la gestion des appels d’outils, de la mémoire et du contexte.}
        \item {Conception d’un pipeline RAG avec chunking, embeddings, indexation dans ChromaDB et recherche sémantique.}
        \item {Développement d’API REST en Python/FastAPI documentées avec OpenAPI, stockage dans MongoDB et interface React/Next.js pour interagir avec l’agent et ses outils.}
        \item {Déploiement de LLM à poids ouverts via vLLM sur un cluster HPC avec GPU NVIDIA et gestion de modèles hébergés localement avec LM Studio sur Mac Studio.}
        \item {Comparaison de LLM et de configurations de quantification avec Slurm, selon le taux de réussite de la recherche, la latence d’inférence et la taille du contexte.}
        \item {Évaluation d’OpenClaw dans un environnement Docker isolé pour étudier les risques liés à l’accès d’un agent aux outils et aux fichiers.}
        \item {Conteneurisation de composants avec Docker et mise en place d’un pipeline CI/CD sur GitHub.}
      \end{cvitems}
    }

  \cventry
    {Stagiaire ingénieur de recherche en apprentissage par renforcement}
    {Université Karunya}
    {Coimbatore, Inde}
    {Juin 2025 -- Septembre 2025}
    {
      \begin{cvitems}
        \item {Développement d’un modèle d’apprentissage par renforcement pour la navigation autonome d’un robot hospitalier, puis déploiement et tests sur Raspberry Pi.}
      \end{cvitems}
    }
\end{cventries}
